\section{Discussion: Policy Design}\label{sec:discussion}

The policy lesson is not that set-asides are universally dominated. It is that bidder-exclusion rules should be evaluated against the thickness of the protected pool, the heterogeneity of the good, and the welfare weight the planner places on SME surplus. Law 14{,}133/2021 allows both full SME set-asides and price preferences. The estimates here imply that the choice between them should be contingent on market structure rather than imposed as a uniform default.

\subsection{When does a set-aside become defensible?}

The structural decomposition suggests a simple answer. A full set-aside becomes more defensible when three conditions line up: the SME pool is thin enough that protecting it is a first-order objective, the good is heterogeneous enough that the identity of the induced entrants matters for equilibrium selection, and the planner places a high welfare weight on SME surplus. Those are precisely the conditions under which the policy ranking becomes model-sensitive in the data.

Non-pharma sits at the opposite corner. The baseline SME pool is thicker, the goods are more standardized, and the \optPrefThreshold{} price preference is Pareto-superior under both the main equilibrium-selection reading and the strict-invariance benchmark. That is the market class in which bidder exclusion is hardest to defend. Pharma is closer to the knife-edge: the pool is thinner, heterogeneity is higher, and the policy ranking depends on whether the post-policy entrant mix is treated as selected equilibrium participation or forced to equal the pre-policy primitive. The point is not that pharma is special. The point is that thin heterogeneous markets are exactly where bidder-exclusion rules should be expected to be more sensitive.

The pharmaceutical bifurcation is the paper's central policy claim, not a hedge against an unwelcome result. The paper does not assert universal preference-rule dominance; it identifies the conditions (thin protected pool, heterogeneous good) under which the standard policy ranking reverses, and it documents that those conditions are met in only one of the two Group-65 sub-classes studied. The pharmaceutical welfare ranking under the main specification depends on interpreting the post-policy SME cost distribution as the result of equilibrium selection rather than as a latent unobserved-heterogeneity shift. Under either reading, the policy ranking is the analytically derivable one given the assumed primitive; the deeper question of which interpretation is empirically correct cannot be settled with a single-reform single-platform dataset. Of the eight (class $\times$ type $\times$ period) cells in the structural sample, the policy reversal between V0 and V3 obtains in exactly the two pharma cells under strict invariance and in zero cells under the main specification; non-pharma is robust to both readings. This is a sharp policy partition, not a fragility that the paper papers over: the conditional reading is the contribution, not a qualifier on a universal contribution.

\subsection{The welfare weight implicit in the current regime}

Following the generalized social-marginal-welfare-weight framework of \citet{saezstantcheva2016}, consider a social planner aggregating producer and consumer surplus with weight $w^{\text{SME}}$ on SME producer surplus. The planner is indifferent between the full set-aside (V0) and the \optPrefThreshold{} price preference (V3) when
\begin{equation}\label{eq:welfare_weight}
w^{\text{SME}}_\star
\;=\; \frac{\mathrm{Loss}(V_0) - \mathrm{Loss}(V_3)}
           {\mathrm{SMESurplus}(V_0) - \mathrm{SMESurplus}(V_3)}.
\end{equation}
The numerator is the loss differential. Under the main specification, $\mathrm{Loss}(V_0) - \mathrm{Loss}(V_3) \approx \welfTotalLossLthirtyNp$ in non-pharma and $\welfTotalLossLthirtyPh$ in pharma: V0 welfare loss at $\lambda = \lambdaMain$ is $\welfTotalLossLthirtyNp$ and $\welfTotalLossLthirtyPh$ (Table~\ref{tab:v3_welfare}); V3 welfare loss is near zero (Table~\ref{tab:v3_apv}, $\Delta = \vThreeDeltaNp$ non-pharma and $\vThreeDeltaPh$ pharma).

The denominator is the SME-surplus differential. Under V0, the implicit transfer to SME winners is $\Delta_{\text{gov}} - \mathrm{DWL}_{\text{alloc}} = \welfTransferNp$ in non-pharma and $\welfTransferPh$ in pharma (Table~\ref{tab:v3_welfare}); under V3 the transfer is approximately zero given the near-zero $\Delta$ on the open-regime auction.\footnote{Under the English-reverse IPV interpretation of Section~\ref{sec:model}, all losers exit at cost and earn zero auction-stage profit; SMESurplus reduces to the winner-conditional rectangular transfer.}

The ratio gives the indifference weight. Substituting in equation~\eqref{eq:welfare_weight}: $w^{\text{SME}}_\star \approx \welfTotalLossLthirtyNp/\welfTransferNp \approx \welfWeightMainNp$ in non-pharma and $\welfTotalLossLthirtyPh/\welfTransferPh \approx \welfWeightMainPh$ in pharma under the main specification. Under strict invariance, non-pharma is unchanged (implied weight still above the utilitarian benchmark and well above the empirical MVPF of Brazil's flagship transfer program), while the pharma indifference weight falls to \welfWeightSiPh (Table~\ref{tab:v3_strict_invariance}). These are not contradictory numbers; they are a map of when bidder exclusion is doing something that a moderate preference rule cannot.

The comparison is informative because it disciplines the distributional argument in favor of set-asides. A planner willing to pay the non-pharma fiscal cost of bidder exclusion would have to value one real of SME surplus at more than twice a real of general welfare, even in a thick market where a preference rule achieves nearly the same allocative objective. In pharma, the required weight depends on how strongly one believes the policy selects a different entrant mix. The right reading is therefore conditional: bidder exclusion is hardest to justify in thick standardized markets, and most plausible only in thin heterogeneous markets with very high welfare weights on protected firms. Figure~\ref{fig:v3_welfare_weight} plots the implied weights against the empirical Bolsa~Fam\'ilia MVPF benchmark of \citet{bergstrom2025}, who estimate the marginal value of public funds of Brazil's flagship cash-transfer program at $\mvpfBolsaLo$--$\mvpfBolsaHi$ depending on whether GE-multiplier effects are included; \citet{mendes2023} document that the GE multiplier is non-trivial in the Brazilian context, consistent with the upper bound of that interval.

\begin{figure}[!htbp]
\centering
\includegraphics[width=0.98\textwidth]{../output/figures/fig_v3_welfare_weight.pdf}
\caption[Implicit welfare weight to prefer the set-aside over the price preference.]{\textit{Implicit welfare weight $w^{\text{SME}}_\star$ required to prefer the full set-aside over a \optPrefThreshold{} price preference, by class and specification.} The grey band marks the empirical MVPF range of $\mvpfBolsaLo$--$\mvpfBolsaHi$ estimated by \citet{bergstrom2025} for the Bolsa~Fam\'ilia program (95\% CI [\mvpfBolsaCILo, \mvpfBolsaCIHi] for the central estimate, [\mvpfBolsaCIUpLo, \mvpfBolsaCIUpHi] for the upper bound that incorporates GE-multiplier effects \citep{mendes2023}); the dashed vertical line marks the utilitarian benchmark $w = 1$. In non-pharmaceuticals, the implied weight sits well above the Bolsa~Fam\'ilia band under the main equilibrium-selection specification ($w^{\text{SME}}_\star = \welfWeightMainNp$); under the strict-invariance benchmark, the preference rule continues to dominate (the implied weight remains above the utilitarian benchmark $w = 1$). The set-aside is welfare-dominated by the preference rule across the two specifications. In pharmaceuticals, the implied weight under the main specification ($w^{\text{SME}}_\star = \welfWeightMainPh$) sits above the band but reverses to \welfWeightSiPh{} under strict invariance, sitting at the lower edge of the Bolsa~Fam\'ilia interval; the pharmaceutical ranking is therefore conditional on how the post-policy SME pool is modeled.}
\label{fig:v3_welfare_weight}
\end{figure}

\subsection{Implementation considerations}

Two features qualify the policy comparison. First, V3 is computed in a Vickrey-equivalent framework; the FPSB equivalent under \citet{maskinriley2000} adjustment is bounded above by approximately \maskinRileyBoundPct~percent of the reference price (Section~\ref{sec:fpsb_bound}), well below the V0--V3 expected-price gap of \vZeroVThreeGapLo--\vZeroVThreeGapHi~percent. The welfare ranking V3 $\succ$ V0 is therefore robust to auction format in both classes; the strict-invariance reversal is a more first-order concern in pharma than the FPSB equilibrium adjustment.

Second, a preference rule is procedurally more complex to implement than a set-aside. The buyer must score bids twice and disclose the winner-selection rule clearly. But this is an operational complication, not a conceptual obstacle. Preference systems have been run elsewhere for decades---the U.S. SBA price-evaluation preference for HUBZone and disadvantaged firms, EU Directive 2014/24 transposition rules in several member states, and the Brazilian federal preference margin under Decree 8.538/2015---and Brazilian procurement law already contemplates them. The relevant policy question is therefore not whether the instrument exists, but in which markets it should replace bidder exclusion.

\subsection{Market thickness and goods heterogeneity}

The non-pharma versus pharma comparison is best read as a statement about the interaction between bidder exclusion and goods characteristics. Pharma auctions have both a thinner SME pool and higher within-class heterogeneity. Once non-SMEs are removed from price formation, those two features raise the expected second-order statistic and make the selected entrant mix more important. Non-pharma auctions are closer to the textbook standardized-good environment in which a moderate preference rule can tilt allocation toward SMEs without materially changing price formation. Operationally, the bifurcation in the data here positions the two cells on opposite sides of a \smesPerAucBandLo--\smesPerAucBandHi{} SMEs-per-auction band: non-pharma at \bneNsmePostNp{} SMEs per auction sits at the upper end and the V3 $>$ V0 ranking holds across both specifications; pharma at \bneNsmePostPh{} SMEs per auction sits at the lower end and the ranking reverses under strict invariance. With only two cell observations, this is a positioning of the two empirical cells relative to each other rather than a precisely identified threshold; broader generalization requires the multi-platform extension flagged below.

The pattern is consistent with the Group-65 split but the paper does not estimate it as an external-validity claim outside the BEC platform. Within structurally similar Brazilian procurement settings, goods with deeper supplier bases and tighter specification standardization are plausible candidates for preference rules rather than exclusion; goods with thinner protected pools, more heterogeneous production technologies, and stronger administrative reasons to value SME surplus are the cases in which the policy comparison becomes more model-sensitive. Comparable extension settings include federal Brazilian procurement on ComprasNet (a larger universe with the same Law 14{,}133 framework), municipal procurement under similar SME defaults, and non-Brazilian SME-set-aside reforms such as the U.S. SBA 8(a) program changes or EU Directive 2014/24 transposition events, where the same conditional-on-thickness mechanism could be tested against alternative institutional baselines. The contribution of the paper is to put that intuition on structural objects in one specific reform: the size of the within-auction component, the magnitude of the participation-margin offset, and the welfare weight required to rationalize the policy under each specification.

Two further structural questions---the IPV restriction in pharmaceutical procurement and the possibility of confounding regulatory events at the March 2018 cutoff---are taken up in Appendices~\ref{app:robustness_ipv} and~\ref{app:cutoff} respectively. Both qualify the pharmaceutical welfare-dominance result rather than the non-pharmaceutical headline; the pharmaceutical claim is conditional on the equilibrium-selection treatment of $F_c^{\text{SME,Post}}$, while the non-pharmaceutical claim is preserved across the IPV-relaxed bounds and the strict-invariance specification.

\subsection{Extensions the paper does not attempt}

Three extensions represent natural follow-on work. (i) Quantity response: the buyer's adjustment to higher prices, which the static model treats as inelastic at observed values. Under a low-elasticity assumption $\varepsilon \in [\epsLow, \epsHigh]$ consistent with essential public-health goods, the quantity adjustment to the simulated \didUpperPct~percent log-price increase would shave the fiscal cost by $\varepsilon \cdot \didUpperDec$, or roughly \ppShavedLo--\ppShavedHi{} percentage points of $p_{S_1}$, and would slightly reduce the allocative DWL component. (ii) Dynamic capacity accumulation: SME firms protected by the rule may build production capacity that persists beyond the policy window, an externally validated mechanism in \citet{ferraz2016} and \citet{szerman2023}. (iii) General-equilibrium SME surplus: the partial-equilibrium SMESurplus measure used in the welfare-weight identity (Section~\ref{sec:discussion}) could be extended to a general-equilibrium framework with option value, capacity persistence, and private-sector spillovers. Each is a margin a fully dynamic-general-equilibrium model would price. The direction of the welfare-weight shift is not pinned down by the current paper: each extension could enlarge the SME-side benefit ex post, and the welfare-weight comparison could move in either direction depending on the dynamic offsets to allocative DWL. None changes the organizing lesson: the welfare ranking of bidder-exclusion rules is conditional on market structure.
